\section{Explicit gauges and exact arithmetic}%
\label{sec:asymex_explicit_gauge}

The compression datum of Theorem~\ref{thm:asym_multiblock} carries its
change of bond coordinates as an invertible map onto a graded coordinate
space, one summand per slot. For each worked example below, that map is an explicit
invertible matrix together with an explicit labelling of the bond
coordinates by slots and positions inside slots.

\begin{definition}[Explicit gauge and entrywise integer coercion]
    \label{def:asymex_explicit_gauge}
    \lean{MPSTensor.complexOfInt, MPSTensor.mulIntTensor, MPSTensor.gaugeOfMatrix}
    \leanok
    \uses{def:mps_tensor, def:mpdo_operator_product}
    Let $\tau$ be a bijection between the graded coordinate space of a slot
    family and the bond index set $\{0,\ldots,D_B-1\}$, and let $G$ and
    $G^{-1}$ be mutually inverse matrices in $\MN{D_B}$. The \emph{explicit
        gauge} attached to $(\tau,G)$ is the invertible map sending a bond
    vector $v$ to the graded vector $b\mapsto(Gv)(\tau b)$. For an integer
    matrix $X$, its \emph{entrywise coercion} is the complex matrix with the
    same entries; for integer matrix product operator tensors $M$ and $N$,
    their integer bond-space product is
    $(M\mathbin{\cdot}N)^{ik}=\sum_jM^{ij}\otimes N^{jk}$.
\end{definition}

\begin{theorem}[Conjugation, coercion and the compression pairs of a slot]
    \label{thm:asymex_explicit_gauge}
    % Principal declaration: MPSTensor.conjMatrix_gaugeOfMatrix.
    \lean{MPSTensor.conjMatrix_gaugeOfMatrix, MPSTensor.complexOfInt_mul,
        MPSTensor.mulTensor_complexOfInt, MPSTensor.mulTensor_smul_complexOfInt,
        MPSTensor.MultiBlockCompression.left_gaugeOfMatrix,
        MPSTensor.MultiBlockCompression.right_gaugeOfMatrix}
    \leanok
    \uses{def:asymex_explicit_gauge, thm:asym_multiblock}
    Conjugation by the explicit gauge attached to $(\tau,G)$ sends a bond
    matrix $A$ to the matrix $GAG^{-1}$ relabelled along $\tau$. The
    entrywise coercion is multiplicative, and the bond-space product of two
    integer tensors each rescaled by one scalar $c$ is $c^2$ times the
    entrywise coercion of their integer bond-space product. If a compression
    datum uses the explicit gauge attached to $(\tau,G)$, then the
    compression map $W_s$ out of the bond space consists of the rows of $G$
    at the coordinates of the slot $s$, and the compression map $V_s$ into
    the bond space consists of the columns of $G^{-1}$ at those same
    coordinates.
\end{theorem}

\begin{proof}
    \leanok
    The matrix of the explicit gauge is $G$ relabelled along $\tau$ in its
    rows, and the matrix of its inverse is $G^{-1}$ relabelled along $\tau$
    in its columns, so conjugation is the stated matrix conjugation followed
    by the relabelling. Multiplicativity of the coercion is the fact that a
    sum of products of integers is the same computed in $\Z$ or in $\C$, and
    the statement about the bond-space product follows by expanding both
    sides entrywise. The two compression maps are the composites of the
    coordinate embedding and projection of the slot with the gauge and its
    inverse; evaluating the matrix products at a single graded coordinate
    leaves one row of $G$ and one column of $G^{-1}$.
\end{proof}

\begin{definition}[One-slot compression data]%
    \label{def:asymex_one_slot}
    \lean{MPSTensor.stackedInt, MPSTensor.oneSlot, MPSTensor.oneSlotMem,
        MPSTensor.MultiBlockCompression.ofConjMatrix,
        MPSTensor.MultiBlockCompression.ofConjInt,
        MPSTensor.MultiBlockCompression.ofEq,
        MPSTensor.MultiBlockCompression.ofScalarFlagFour}
    \leanok
    \uses{def:asymex_explicit_gauge, thm:asym_multiblock}
    For a stacked product of two integer matrix product operator tensors,
    read over the pair alphabet, the letter at $a=do+i$ is
    $(M\mathbin{\cdot}N)^{oi}$. A \emph{one-slot compression datum} for a
    tensor $B$ onto a single target $C$ is the data of
    Theorem~\ref{thm:asym_multiblock} for the one-element slot set: an
    explicit gauge $(\tau,G)$ conjugating $B$ into a matrix that is block
    upper triangular for a flag with $C$ at one diagonal position. Two
    shapes recur in the examples below: a conjugation $GB^iG^{-1}=C^i$ for
    an invertible $G$ gives such a datum with no zero slots ($z=0$); and,
    for a four-dimensional $B$ and a one-dimensional $C$, an integer gauge
    sending every letter to a matrix with $C$ at the third diagonal
    position of a four-step flag and zero at the other three gives such a
    datum with three zero slots ($z=3$).
\end{definition}

\begin{theorem}[The remainder of a conjugate datum vanishes; word traces
        of a one-slot datum]%
    \label{thm:asymex_one_slot}
    \lean{MPSTensor.MultiBlockCompression.remainder_ofConjMatrix,
        MPSTensor.MultiBlockCompression.remainder_ofEq,
        MPSTensor.MultiBlockCompression.remainder_ofConjInt,
        MPSTensor.MultiBlockCompression.trace_evalWord_oneSlot,
        MPSTensor.MultiBlockCompression.trace_evalWord_oneSlot_smul}
    \leanok
    \uses{def:asymex_one_slot}
    The remainder of a conjugate one-slot datum vanishes identically: the
    compression is an exact change of coordinates. For every one-slot
    datum, the trace of every nonempty word of $B$ equals the trace of the
    corresponding word of $C$, or, if the target is carried with a scalar
    weight $c$, that power of $c$ times the trace of the word of the
    unweighted target.
\end{theorem}

\begin{proof}
    \leanok
    \uses{thm:asymex_explicit_gauge, prop:asym_compression_trace}
    Immediate from Theorem~\ref{thm:asymex_explicit_gauge} and
    Proposition~\ref{prop:asym_compression_trace}, specialized to a single
    slot.
\end{proof}

\begin{theorem}[Integer certificates for normality and for vanishing
        sitewise intertwiners]%
    \label{thm:asymex_int_certificates}
    \lean{MPSTensor.isNBlkInjective_two_of_int,
        MPSTensor.right_intertwiner_eq_zero_of_int,
        MPSTensor.left_intertwiner_eq_zero_of_int}
    \leanok
    \uses{def:normal}
    If every matrix unit is, up to a common nonzero integer factor, an
    integer combination of length-two words of an integer tensor, the
    length-two words span the full matrix algebra. If the rows,
    respectively columns, of the letters occurring in a sitewise
    intertwiner system form an integer matrix with an integer left inverse
    up to a nonzero factor, the only solution of that system is zero.
\end{theorem}

\begin{proof}
    \leanok
    Clearing the nonzero integer factor by scalar division after coercion
    to $\C$ reduces each statement to the corresponding identity between
    integer matrices, a finite verification.
\end{proof}
